% ============================================================================
% Results
% ============================================================================
\section{Results}
\label{sec:results}

We present our reproduction of the key findings from the original study.

\subsection{Metastable Dynamics in Clustered Networks}
\label{sec:metastable}

An excitatory-inhibitory (E/I) network with clustered E$\to$E connectivity
produces metastable ``high'' firing activity states within clusters that are
stable yet quick to transition. These high activity states persist for
hundreds of milliseconds to seconds before switching, leading to slow rate
dynamics and elevated trial-to-trial spike count variability (super-Poisson)
compared to unclustered networks.

\begin{figure}[t]
    \centering
    \includegraphics[width=\columnwidth]{rasters_uniform}
    \caption{Spike rasters from a uniform (unclustered) network across
    multiple trials.}
    \label{fig:rasters_uniform}
\end{figure}

\begin{figure}[t]
    \centering
    \includegraphics[width=\columnwidth]{rasters_clustered}
    \caption{Spike rasters from a clustered network across multiple trials.}
    \label{fig:rasters_clustered}
\end{figure}

\FloatBarrier
\subsubsection{Mean Rate}
\label{sec:mean_rate}

Despite the switching between high and low firing states, the clustered
network exhibits a low mean firing rate similar to that of the unstructured
network (Figure~\ref{fig:firing_rates}). This is likely due to additional
suppression within the network during high firing activity states. The high
activity epochs in clustered networks do not alter the overall mean firing
rate; however, the firing rate distribution broadens somewhat, likely
reflecting the alternation between high and low activity states.

\begin{figure}[t]
    \centering
    \includegraphics[width=\columnwidth]{firing_rate_distribution}
    \caption{Distribution of firing rates for excitatory neurons in uniform
    and clustered networks. Both networks maintain low mean firing rates,
    though the clustered network shows a broader distribution reflecting
    the alternation between high and low activity states. Triangles indicate
    mean values.}
    \label{fig:firing_rates}
\end{figure}

\subsubsection{Variability}
\label{sec:variability}

Neurons in the clustered network exhibit Fano factors greater than one
(``super-Poisson'' spike count variability), as shown in
Figure~\ref{fig:fano_distribution}. Furthermore, the Fano factor increases
with the size of the counting window (Figure~\ref{fig:fano_vs_window}). This
pattern illustrates the presence of slow firing rate variability in clustered
networks---larger windows capture more of the slow rate fluctuations.

\begin{figure}[t]
    \centering
    \includegraphics[width=\columnwidth]{fano_factor_distribution}
    \caption{Distribution of Fano factors across neurons for uniform and
    clustered networks. The clustered network exhibits a broader distribution
    with a higher mean (triangles), indicating greater spike count
    variability.}
    \label{fig:fano_distribution}
\end{figure}

\begin{figure}[t]
    \centering
    \includegraphics[width=\columnwidth]{fano_factor_vs_window}
    \caption{Mean Fano factor as a function of counting window size. In the
    clustered network, the Fano factor increases with window size, reflecting
    slow rate fluctuations. The uniform network remains near the Poisson
    baseline (FF~=~1).}
    \label{fig:fano_vs_window}
\end{figure}

\FloatBarrier
\subsubsection{Correlations}
\label{sec:correlations}

Pairwise spike count correlations remain near zero on average for both
network types (Figure~\ref{fig:correlations_all}). However, in the clustered
network, within-cluster neuron pairs show a heavier positive tail in the
distribution of correlation coefficients (Figure~\ref{fig:correlations_cluster}).
This observation is consistent with the notion that clusters undergo
additional correlated firing rate transitions---switching between high and
low activity states---when examined across temporal windows spanning the
simulation. Furthermore, this positive tail becomes increasingly pronounced
as the counting window size increases (50~ms vs.\ 100~ms), as expected given
that within-cluster neurons share these slow rate fluctuations.

\begin{figure}[t]
    \centering
    \includegraphics[width=\columnwidth]{correlation_all_pairs}
    \caption{Distribution of pairwise spike count correlations for all
    neuron pairs. Correlations remain centered near zero on average for
    both network types.}
    \label{fig:correlations_all}
\end{figure}

\begin{figure}[t]
    \centering
    \includegraphics[width=\columnwidth]{correlation_same_cluster}
    \caption{Distribution of pairwise spike count correlations for
    within-cluster pairs only. The clustered network shows a heavy positive
    tail that increases with counting window size (50~ms vs.\ 100~ms).}
    \label{fig:correlations_cluster}
\end{figure}

\subsubsection{Timescales}
\label{sec:timescales}

Clustered networks exhibit substantially longer correlation times than
unstructured networks. The autocovariance function decays much more slowly
as a function of lag in clustered networks compared to their unstructured
counterparts. Similarly, the cross-covariance for within-cluster neuron
pairs decays slowly, indicating a shared correlated rate change among
neurons belonging to the same cluster.

% ============================================================================
% Critical Cluster Strength
% ============================================================================
\FloatBarrier
\subsection{Critical Cluster Strength}
\label{sec:critical_ree}

TODO

\begin{figure}[t]
    \centering
    \includegraphics[width=\columnwidth]{fano_factor_vs_ree}
    \caption{Mean Fano factor as a function of cluster strength ($R_{EE}$).
    As within-cluster connection strength increases, the network transitions
    from Poisson-like variability to super-Poisson variability. The dashed
    line indicates the Poisson baseline (FF~=~1).}
    \label{fig:fano_vs_ree}
\end{figure}
